%%% General macro and environment setup %%%

% General typesetting %
% Simultaneous sub- and superscript in text
% Taken from https://tex.stackexchange.com/a/24082/414595
\makeatletter
\DeclareRobustCommand*\textsubsuperscript[2]{%
  \@textsubsuperscript{\selectfont#1}{\selectfont#2}}
\def\@textsubsuperscript#1#2{%
  {\m@th\ensuremath{_{\mbox{\fontsize\sf@size\z@#1}}
                    ^{\mbox{\fontsize\sf@size\z@#2}}}}}
\makeatother


% Boxes %
% Boxes with dashed outlines
%% Default distance between content and outline in tight boxes
\newlength{\tboxsep}
\setlength{\tboxsep}{0.5pt}

%% Context-aware (math) box with regular/solid outlines
\newcommand{\boxedmp}[1]{\mathpalette{\boxedmpaux}{#1}}
\newcommand{\boxedmpaux}[2]{\boxed{#1#2}}

%% Tight solid box (text mode)
\NewDocumentCommand{\tbox}{o m}{
  \IfNoValueTF{#1}{%
    \setlength{\fboxsep}{\tboxsep}%
  }{%
    \setlength{\fboxsep}{#1}%
  }%
  \fbox{#2}%
}

%% Tight solid box (math mode)
\NewDocumentCommand{\tboxed}{o m}{
  \IfNoValueTF{#1}{%
    \setlength{\fboxsep}{\tboxsep}%
  }{%
    \setlength{\fboxsep}{#1}%
  }%
  \boxedmp{#2}%
}

%% Context-aware (math) box with dashed outlines
\newcommand{\dboxmp}[1]{\mathpalette{\dboxmpaux}{#1}}
\newcommand{\dboxmpaux}[2]{\dbox{\ensuremath{#1#2}}}

%% Tight dashed box (text mode)
\NewDocumentCommand{\tdbox}{o m}{
  \IfNoValueTF{#1}{%
    \setlength{\fboxsep}{\tboxsep}%
  }{%
    \setlength{\fboxsep}{#1}%
  }%
  \dbox{#2}%
}

%% Tight dashed box (math mode)
\NewDocumentCommand{\tdboxed}{o m}{
  \IfNoValueTF{#1}{%
    \setlength{\fboxsep}{\tboxsep}%
  }{%
    \setlength{\fboxsep}{#1}%
  }%
  \dboxmp{#2}%
}

%% Tight dotted box (text mode)
\NewDocumentCommand{\tdtbox}{o m}{
  \IfNoValueTF{#1}{%
    \setlength{\fboxsep}{\tboxsep}%
  }{%
    \setlength{\fboxsep}{#1}%
  }%
  \setlength{\dashlength}{2pt}%
  \setlength{\dashdash}{0.5pt}%
  \dbox{#2}%
}

%% Tight dotted box (math mode)
\NewDocumentCommand{\tdtboxed}{o m}{
  \IfNoValueTF{#1}{%
    \setlength{\fboxsep}{\tboxsep}%
  }{%
    \setlength{\fboxsep}{#1}%
  }%
  \setlength{\dashlength}{2pt}%
  \setlength{\dashdash}{0.5pt}%
  \dboxmp{#2}%
}


% Mathematics %
% Theorem, definitions, and remarks
%% Theorem-like
\declaretheoremstyle[%
headfont=\normalfont\normalsize\normalcolor\bfseries,%
notefont=\normalfont\normalsize\normalcolor\mdseries,%
bodyfont=\normalfont\normalsize\normalcolor\slshape,%
postheadspace=0.5em,%
]{thm}

\declaretheorem[%
style=thm,%
name=Theorem,%
% refname={theorem,theorems},%
% Refname={Theorem,Theorems},%
]{theorem}

\declaretheorem[%
style=thm,%
name=Lemma,%
% refname={lemma,lemmas},%
% Refname={Lemma,Lemmas},%
numberlike=theorem,%
]{lemma}

\declaretheorem[%
style=thm,%
name=Axiom,%
% refname={axiom,axioms},%
% Refname={Axiom,Axioms},%
numberlike=theorem,%
]{axiom}

\declaretheorem[%
style=thm,%
name=Corollary,%
% refname={corollary,corollarys},%
% Refname={Corollary,Corollarys},%
numberlike=theorem,%
]{corollary}

\declaretheorem[%
style=thm,%
name=Proposition,%
% refname={proposition,propositions},%
% Refname={Proposition,Propositions},%
numberlike=theorem,%
]{proposition}

%% Definition-like
\declaretheoremstyle[%
headfont=\normalfont\normalsize\normalcolor\bfseries,%
notefont=\normalfont\normalsize\normalcolor\mdseries,%
bodyfont=\normalfont\normalsize\normalcolor,%
postheadspace=0.5em,%
]{def}

\declaretheorem[%
style=def,%
name=Definition,%
% refname={definition,definitions},%
% Refname={Definition,Definitions},%
]{definition}

%% Remark-like
\declaretheoremstyle[%
headfont=\normalfont\normalsize\normalcolor\itshape,%
notefont=\normalfont\normalsize\normalcolor,%
bodyfont=\normalfont\normalsize\normalcolor,%
postheadspace=0.5em,%
]{rmrk}

\declaretheorem[%
style=rmrk,%
name=Remark,%
% refname={remark,remarks},%
% Refname={Remark,Remarks},%
]{remark}

\declaretheorem[%
style=rmrk,%
name=Claim,%
% refname={claim,claims},%
% Refname={Claim,Claims},%
]{claim}


% Code %
% Sampling
% Based on implementation from cryptocode package
% (see https://ctan.org/pkg/cryptocode and https://tex.stackexchange.com/a/418791/414595)
\makeatletter
\newcommand{\smalldollar}{\mathrel{\mathpalette\small@dollar\relax}}
\newcommand{\small@dollar}[2]{%
  \vcenter{\hbox{%
    $#1\textnormal{\fontsize{0.7\dimexpr\f@size pt}{0}\selectfont\$}$%
  }}%
}
\makeatother

\newcommand*{\sample}{\leftarrow\mathrel{\mkern-1.5mu}\smalldollar}
